\chapter{Arquitectura del Sistema}
\label{cap:arquitectura}

El diseño lógico de este proyecto orquesta una extensa pila tecnológica (\textit{Stack}) desacoplada para garantizar escalabilidad, monitorización exhaustiva y una separación de responsabilidades estricta. La arquitectura general emplea una metodología \textit{microservices-oriented} (orientada a microservicios) empaquetada globalmente mediante contenedores Docker con un manifiesto \textit{Docker Compose}.

% ============================================================
\section{Vista General en Capas}
% ============================================================
El sistema se diseña partiendo de una distinción clásica entre la fachada de interfaz visible, la lógica de control intermedia y múltiples bases de datos de apoyo especializadas. Se pueden identificar cinco capas claramente diferenciadas, desde la interfaz hasta la persistencia.

\subsection{Capa de Presentación (Frontend)}
El único punto de acceso de cara al usuario está gobernado por el módulo OpenWebUI e iterado a través de un proxy inverso Nginx, asegurando una entrada HTTPS fortificada con puertos redireccionados aislados. Esta capa contiene el componente web que implementa la pasarela para interaccionar nativamente con el motor, además de gestionar la Autenticación Única (SSO, \textit{Single Sign-On}) del Azure Active Directory.

OpenWebUI fue seleccionado frente a alternativas como Gradio, Streamlit o ChatUI por su compatibilidad nativa con Ollama, su sistema de plugins (\textit{Pipelines}) que permite inyectar lógica personalizada sin modificar el código fuente del frontend, y su gestión de usuarios y conversaciones integrada de serie con persistencia en PostgreSQL.

\subsection{Capa de Seguridad Perimetral}
El proxy inverso NGINX actúa como la primera barrera de seguridad del ecosistema. Escucha exclusivamente en el puerto \textbf{8443} con un certificado TLS/SSL autofirmado para el tráfico interno corporativo. Su configuración principal incluye:

\begin{itemize}
    \item \textbf{Terminación TLS:} Todo el tráfico entre el navegador del usuario y el servidor se cifra. Los microservicios internos se comunican en texto plano dentro de la red Docker aislada.
    \item \textbf{Proxy pass transparente:} Las peticiones HTTPS se redireccionan internamente a OpenWebUI (puerto 8080) preservando las cabeceras \texttt{X-Forwarded-For} y \texttt{X-Real-IP} para la trazabilidad.
    \item \textbf{WebSocket upgrade:} Se habilita la mejora de protocolo a WebSocket, requisito indispensable para el streaming de tokens en tiempo real (\textit{Server-Sent Events}).
\end{itemize}

\subsection{Capa de Orquestación y Enrutamiento Inteligente}
De manera subyacente a OpenWebUI, un agente enrutador denominado \textit{JARVIS Pipeline} recibe cada evento. Elaborado sobre Python nativo, JARVIS no emite inferencias instantáneas, sino que evalúa la intención del usuario a nivel semántico (por ejemplo, detectar la mención de hipervínculos para un \textit{scraping} asíncrono, deducir la motivación de emplear el corpus documental mediante la activación del \textit{RAG} explícito, o dictaminar la necesidad de contactar al BOE). JARVIS orquesta flujos mediante peticiones asíncronas HTTP a la capa de Aplicación según los parámetros resueltos.

El pipeline opera como un \textit{middleware} interpuesto entre la interfaz de usuario y los modelos de lenguaje. OpenWebUI envía a JARVIS cada mensaje del usuario como un evento HTTP con el texto, los adjuntos (si los hay) y los metadatos del usuario. JARVIS analiza el contenido, selecciona el modo de operación apropiado y delega la ejecución al microservicio correspondiente.

\subsection{Capa de Aplicación (Backend)}
El cerebro lógico principal del proyecto se sustenta en tres microservicios interconectados:

\begin{itemize}
    \item \textbf{API Backend (FastAPI):} Servicio central que atiende todas las operaciones RAG. Implementa los endpoints REST para la ingesta de documentos (\texttt{/ingest}), la búsqueda semántica (\texttt{/chat}), el web scraping (\texttt{/scrape}), la consulta de documentos indexados (\texttt{/documents}) y la exportación de métricas para Prometheus (\texttt{/metrics}). Internamente orquesta la cadena de procesamiento: chunking semántico, generación de embeddings con SentenceTransformer, almacenamiento en Qdrant y formateo de prompts contextuales para el LLM.

    \item \textbf{Indexer API:} Responsable explícito de sondear periódicamente los orígenes asíncronos. Realiza la conexión a \textit{Microsoft Graph} monitorizando novedades en los ficheros de SharePoint o las peticiones en carpetas locales. Llama automáticamente al motor OCR (PaddleOCR) habilitado sobre GPU si detecta PDFs escaneados a gran escala o ficheros imposibles de parsear semánticamente por texto plano. Funciona de forma completamente autónoma mediante un cron interno configurable.

    \item \textbf{LiteLLM Proxy:} Intermediario en el control de acceso a los Modelos Locales, unificando endpoints de Ollama bajo una API compatible con el estándar de OpenAI. Registra tiempos de servicio, habilita cachés perimetrales sobre Redis para evitar múltiples llamadas de red de las respuestas computadas y aplica reglas de balanceo de carga cuando varios modelos están disponibles simultáneamente.
\end{itemize}

\subsection{Capa de Datos}
La persistencia de memoria se divide drásticamente en un entorno distribuido en cuatro nodos independientes en la pila:

\begin{enumerate}
    \item \textbf{Qdrant --- Base de Datos Vectorial:} Administra las colecciones \texttt{documents} y \texttt{webs} almacenando incrustaciones generadas (\textit{Embeddings} a 384 dimensiones) indexadas por índices hiper-dimensionales afines HNSW (\textit{Hierarchical Navigable Small World}). Despliega además en conjunción estrategias de filtrado disperso como BM25 para búsqueda híbrida. La colección de documentos se segmenta por departamento, permitiendo el control de acceso granular.

    \item \textbf{PostgreSQL con pgvector:} Garantiza consistencia relacional ACID en cuanto a persistir el árbol histórico de interacciones (conversaciones, turnos de palabras, usuarios registrados y trazabilidad del SSO). La extensión pgvector se mantiene instalada como capacidad futura para operaciones vectoriales secundarias.

    \item \textbf{Redis --- Almacén en Memoria:} Interviene como mecanismo de caché para aligerar peticiones al LLM de respuestas repetitivas o idénticas. Cada respuesta cacheada tiene un \textit{Time To Live} (TTL) configurable de 1 hora, tras el cual se invalida automáticamente.

    \item \textbf{MinIO S3 --- Almacenamiento de Objetos:} Preserva una réplica aislada de todos los objetos en formato archivo binario en el sistema, dotando de resistencia a la configuración frente a fallos y habilitando un sistema de backup independiente de la base de datos.
\end{enumerate}

% ============================================================
\section{Orquestación con Docker Compose}
% ============================================================

La totalidad del ecosistema se define declarativamente en un único fichero \texttt{docker-compose.yml}. Este manifiesto YAML describe los 12 servicios del sistema, sus dependencias, volúmenes persistentes, variables de entorno y políticas de reinicio.

\begin{table}[H]
\centering
\caption{Servicios definidos en el manifiesto Docker Compose.}
\label{tab:docker-services}
\begin{tabular}{llc}
\toprule
\textbf{Servicio} & \textbf{Imagen / Build} & \textbf{Puerto Expuesto} \\
\midrule
nginx           & nginx:alpine             & 8443 (HTTPS) \\
open-webui      & open-webui:main          & 8080 (interno) \\
pipelines       & Custom (Dockerfile)      & 9099 (interno) \\
rag-backend     & Custom (Dockerfile)      & 8000 (interno) \\
indexer         & Custom (Dockerfile)      & 8001 (interno) \\
ollama          & ollama/ollama:latest     & 11434 (interno) \\
litellm         & litellm/litellm:main     & 4000 (interno) \\
qdrant          & qdrant/qdrant            & 6333 (interno) \\
postgres        & postgres:16              & 5432 (interno) \\
redis           & redis:7-alpine           & 6379 (interno) \\
mcp-boe-server  & Custom (Dockerfile)      & 8010 (interno) \\
minio           & minio/minio              & 9000 (interno) \\
\bottomrule
\end{tabular}
\end{table}

Todos los servicios comparten una red Docker interna de tipo \textit{bridge} denominada \texttt{tfg-network}. Únicamente el servicio de Nginx expone un puerto al exterior (8443), garantizando que ningún microservicio interno sea accesible directamente desde fuera del servidor.

Los servicios que requieren acceso a GPU (Ollama, PaddleOCR) utilizan la directiva \texttt{deploy.resources.reservations.devices} de Docker Compose para reservar dispositivos NVIDIA mediante el runtime \texttt{nvidia-container-toolkit}.

% ============================================================
\section{Modelos de Inteligencia Artificial Desplegados}
% ============================================================

El sistema opera simultáneamente con seis modelos de IA ejecutados íntegramente en hardware local. La selección de modelos se realizó en base a un equilibrio entre capacidad generativa, consumo de VRAM y velocidad de inferencia.

\begin{table}[H]
\centering
\caption{Inventario completo de modelos IA desplegados y su configuración en LiteLLM.}
\label{tab:modelos-litellm}
\begin{tabular}{lllc}
\toprule
\textbf{Alias LiteLLM} & \textbf{Modelo Real} & \textbf{Función} & \textbf{VRAM} \\
\midrule
\texttt{llama3.1-8b}      & llama3.1:8b-instruct-q8\_0     & Chat y RAG general  & $\sim$8 GB \\
\texttt{llama3.1-summarize}& llama3.1:8b-instruct-q4\_0     & Resúmenes ligeros   & $\sim$4 GB \\
\texttt{JARVIS}            & rag-qwen-ft:latest             & RAG corporativo (FT)& $\sim$5 GB \\
---                        & qwen2.5vl:7b                   & Visión multimodal   & $\sim$5 GB \\
---                        & llava:13b                      & Visión (legacy)     & $\sim$8 GB \\
---                        & MiniLM-L12-v2                  & Embeddings (CPU)    & 0 GB \\
\bottomrule
\end{tabular}
\end{table}

La cuantización Q8\_0 del modelo LLaMA 3.1 8B ofrece un compromiso óptimo entre calidad y velocidad: reduce el consumo de VRAM un 50\% respecto a la versión FP16 con una degradación de calidad inferior al 1\% en benchmarks estándar. El modelo Q4\_0 se reserva para tareas de resumen donde la fidelidad absoluta es menos crítica.

% ============================================================
\section{Flujo de Datos End-to-End}
% ============================================================

Para ilustrar la cohesión del sistema, se describe el flujo completo que atraviesa una consulta RAG típica desde que el usuario la formula hasta que recibe la respuesta:

\begin{enumerate}
    \item El \textbf{usuario} escribe su pregunta en OpenWebUI y pulsa Enter.
    \item \textbf{NGINX} recibe la petición HTTPS y la redirecciona internamente a OpenWebUI.
    \item \textbf{OpenWebUI} envía el mensaje al pipeline JARVIS mediante una llamada HTTP POST al puerto 9099.
    \item \textbf{JARVIS} analiza el texto, detecta la intención RAG y construye una petición al Backend con el query del usuario y las colecciones autorizadas para su departamento.
    \item El \textbf{Backend} genera el embedding de la consulta con SentenceTransformer, ejecuta una búsqueda híbrida (HNSW + BM25) en Qdrant y recupera los $k$ fragmentos más relevantes.
    \item El Backend compone un \textit{prompt} estructurado que incluye el contexto recuperado, la pregunta original y las instrucciones de sistema, y lo envía a LiteLLM.
    \item \textbf{LiteLLM} verifica si existe una respuesta cacheada en Redis. Si no, delega al modelo apropiado en Ollama.
    \item \textbf{Ollama} ejecuta la inferencia en GPU y devuelve los tokens generados en streaming.
    \item La respuesta fluye de vuelta por la cadena: LiteLLM $\rightarrow$ Backend $\rightarrow$ JARVIS $\rightarrow$ OpenWebUI $\rightarrow$ NGINX $\rightarrow$ navegador del usuario.
\end{enumerate}

Todo el flujo se completa típicamente en menos de 3.5 segundos para consultas RAG estándar, con el primer token visible para el usuario en menos de 1.2 segundos gracias al streaming.

% ============================================================
\section{Redes, Seguridad y Monitorización}
% ============================================================

\subsection{Aislamiento de Red}
Toda la malla de nodos convive dentro de un plano interno \texttt{tfg-network}, un modo de puente o \textit{bridge network} gestionado por el motor host de Docker, impidiendo a entes externos o atacantes conectar con puertos de base de datos no expuestos frontalmente. Esto significa que servicios como PostgreSQL, Redis y Qdrant \textbf{no son accesibles} desde fuera del servidor bajo ninguna circunstancia.

\subsection{Autenticación SSO con Azure AD}
OpenWebUI se configura con el protocolo OpenID Connect (OIDC) apuntando al \textit{tenant} corporativo de Azure Active Directory. Los usuarios se autentican con sus credenciales corporativas sin necesidad de crear una cuenta separada. Los tokens JWT emitidos por Azure AD se validan en cada petición y sus atributos (nombre, correo, grupos) se utilizan para el control de acceso departamental en JARVIS.

\subsection{Stack de Observabilidad}
Para consolidar la resiliencia en un marco empresarial de alto rendimiento, se integran los siguientes componentes de monitorización:

\begin{itemize}
    \item \textbf{Prometheus:} Servidor de métricas temporales que scrapea activamente los endpoints \texttt{/metrics} del Backend (FastAPI) cada 15 segundos, almacenando series temporales de latencia, throughput, errores y uso de recursos.

    \item \textbf{NVIDIA DCGM Exporter:} Exportador especializado que captura la saturación paralela de las GPUs de NVIDIA: consumo de memoria GDDR6 (VRAM), utilización de los \textit{CUDA cores}, temperatura y potencia. Estas métricas son scrapeadas por Prometheus y resultan críticas para detectar cuellos de botella en la inferencia.

    \item \textbf{Grafana:} Portal de visualización que consume las métricas de Prometheus y las presenta en dashboards interactivos. Se configuraron tres dashboards principales: uno para métricas del Backend (latencia, peticiones, errores), otro para métricas de GPU (VRAM, temperatura, carga) y un tercero para métricas de base de datos.

    \item \textbf{pgAdmin 4:} Interfaz web de administración de PostgreSQL que permite auditar directamente las tablas de conversaciones, usuarios y sesiones almacenadas por OpenWebUI.
\end{itemize}
